\section{Block fitting and misclassification control}
\label{restore:fit}

Write $\Pi=UU^\top$. For the
proof, rescale the encoded adjacency matrix by its common threshold:
$Q=\widetilde A/\tau=G^{-1/2}AG^{-1/2}$. This changes the eigenvalue scale,
not the retained eigenspace. We work on the aggregate event
\begin{align}
 \norm{Q-G^{-1/2}{ZPZ^\top} G^{-1/2}}_{\op}
   &\le C/\sqrt {np},& \sum_iG_i&\le C{n^2p},
 \label{restore:fit-geometry-one}\\
 \sum_i\frac{(G_i-\tau)^2+(D_i-\sum_jP_{z_i z_j})^2}{G_i}
   &\le Cn,&
 \norm{(I-\Pi)G^{1/2}Z}_F&\le C\sqrt n,
 \label{restore:fit-geometry-two}\\
 \max_{k,\ell}\left|[Z^\top(A-{ZPZ^\top})Z]_{k\ell}\right|
   &\le Cn.&&\label{restore:fit-block-event}
\end{align}
We also have $G_i\ge\tau\asymp {np}$, $\sum_jP_{z_i z_j}\asymp {np}$,
$\bar D\asymp {np}$, and $\tau\ge\max_i\sum_jP_{z_i z_j}$. The first two
lines follow from \cref{prop:sparse-floor}; the block-total event is
paid in the risk proof. All bounds in this section are deterministic
on this event. Constants depend only on the fixed structural and
aggregate-event bounds and, for local estimates, on the fixed thresholds
$L_D,L_U$. They do not depend on the fitted partition. As in the
parameter-estimation and classification arguments of \citet{zhouli}, we display the dependence on misclassification rather
than absorbing it into a separately named error-radius constant.
Align the current labels once with the true labels. Then the incorrectly
assigned set $E=\{i:\widehat z_i\ne z_i\}$ has
$|E|=n\Mis(\widehat z,z)$. The iteration argument below preserves this
same label matching.
For current labels $\widehat z$, let
$\widehat Z_{ik}=\ind{\widehat z_i=k}$, and use
$n_k(\widehat z)$ for their community sizes.

An aligned population basis is
\[
 V=G^{-1/2}Z(Z^\top G^{-1}Z)^{-1/2}O,
\]
where $O$ is orthogonal. The signal gap gives
\begin{equation}\label{restore:fit-basis}
 \norm{U-V}_F\le C/\sqrt {np},\qquad
 \sqrt{G_i}\norm{V_{i\cdot}}_2\le C\sqrt{p}.
\end{equation}
Moreover $\norm Q_{\op}\le1$, since
$|x^\top Ax|\le\sum_iD_ix_i^2\le\sum_iG_ix_i^2$.
In particular $\norm{\Lambda/\tau}_{\op}\le1$, and
$\widehat A=G^{1/2}Q_KG^{1/2}$ for $Q_K=U(\Lambda/\tau)U^\top$.

For the row bounds in this section, put
$a_i=\sqrt{G_i}\norm{U_{i\cdot}}_2$.

\begin{lemma}[Weighted mass of incorrectly assigned vertices]
\label{restore:fit-wrong-mass}
For every set $E\subseteq[n]$,
\begin{equation}\label{restore:fit-wrong-mass-bounds}
 \sum_{i\in E}G_i\le C(n+np|E|),\qquad
 \sum_{i\in E}a_i\le C\left(|E|\sqrt p+\frac1{\sqrt p}\right).
\end{equation}
Also $\sum_i a_i\le Cn\sqrt p$. These bounds hold simultaneously
for all sets; in particular, the error set has
$|E|=n\Mis(\widehat z,z)$ under the common label matching.
\end{lemma}

\begin{proof}
The inequality $G_i\le2\tau+(G_i-\tau)^2/G_i$ and the aggregate
energy bound give $\sum_EG_i\le C(np|E|+n)$. By
\eqref{restore:fit-basis} and Cauchy--Schwarz,
\[
 \sum_{i\in E}a_i
 \le C|E|\sqrt p+
      \left(\sum_{i\in E}G_i\right)^{1/2}\norm{U-V}_F
 \le C\left(|E|\sqrt p+\sqrt{|E|+p^{-1}}\right).
\]
The inequality $\sqrt{p|E|+1}\le p|E|+1$ yields the second bound.
The global bound follows from
$\sum_i a_i\le(\sum_iG_i)^{1/2}\norm U_F$ and fixed $K$.
\end{proof}

\begin{lemma}[Radius-count initialization]
\label{restore:fit-initializer}
The initializer in \cref{restore:algorithm} satisfies
\[
 \Mis(\widehat z,z)\le \frac{C}{np}
\]
on the aggregate event. The comparison uses one common label permutation.
\end{lemma}

\begin{proof}
The normalized rows of $V$ are the $K$ rows of $O$, with pairwise
distance $\sqrt2$. Choose a fixed $B_0$ larger than twice the
uniform upper bound on $\tau/(np)$. Each row with $G_i>B_0{np}$
contributes at least $B_0{np}/4$ to $\sum_i(G_i-\tau)^2/G_i$, so
there are at most $\frac{C}{p}$ such rows. For every remaining row,
$\norm{V_{i\cdot}}_2\ge c/\sqrt n$. The normalization inequality
\[
 \left\|\frac u{\norm u_2}-\frac v{\norm v_2}\right\|_2
 \le\frac{2\norm{u-v}_2}{\norm v_2}
\]
(with any fixed unit direction when $u=0$) and
\eqref{restore:fit-basis} imply that the normalized empirical rows
$x_i$ have total squared distance at most $\frac{C}{p}$ from their true
centers $O_{z_i\cdot}$: the remaining rows satisfy this by the displayed
inequality, and the $O(\frac{1}{p})$ high-degree rows each contribute at most four.

Call a row good if $\norm{x_i-O_{z_i\cdot}}_2\le1/32$.
At most $m=\frac{C}{p}$ rows are bad, and $m<n/(3\beta K)$ for large $n$.
Inductively, all good rows of a selected block have been removed and
all good rows of unselected blocks remain. A radius-$1/8$ ball
centered at a good row of an unselected block contains every good row
of that block, hence more than $m$ remaining rows. The maximizing
ball must therefore contain a good row from an unselected block;
its representative is within $1/8+1/32=5/32$ of that center.
Every good row of that block lies within $3/16<1/4$ of the
representative. Good rows of any other block are farther than
$\sqrt2-3/16>1/4$, so deletion preserves the induction.
After $K$ selections, each good row is closer to its own
representative than to all others. Only the $m$ bad rows can be
misclassified. Dividing by $n$ gives the stated misclassification rate.
\end{proof}

For the remaining reconstruction estimates, write
$v=A\one-\widehat A\one$ and $R_1=\widehat A+v\one^\top/n$.
These are proof quantities computed from the same fixed reconstruction;
they are not additional algorithmic updates.

\begin{lemma}[Reconstructed block totals]
\label{restore:fit-block-totals}
The hollow, degree-preserving reconstruction in
\eqref{restore:reconstruction} satisfies
\begin{equation}\label{restore:fit-global-absolute}
 \sum_{i,j}|\widehat A^{\mathrm{c}}_{ij}|\le Cn^2p.
\end{equation}
For every estimated partition, under the common label matching,
\begin{equation}\label{restore:fit-total-accuracy}
 \left|[\widehat Z^\top\widehat A^{\mathrm{c}}\widehat Z]_{k\ell}
                  -n_kn_\ell P_{k\ell}\right|
 \le Cn\{1+np\Mis(\widehat z,z)\}.
\end{equation}
The omitted diagonal is included in the bound, and the constant is
uniform over the estimated partition.
\end{lemma}

\begin{proof}
With $r_i=(R_1)_{ii}$, the reconstruction formula is equivalently
\[
 \widehat A^{\mathrm{c}}=R_1+\frac{r\one^\top}{n-1}
                   -\frac n{n-1}\diag(r),
\]
because its off-diagonal entries are $\widehat A_{ij}+(v_i+\widehat A_{ii})/(n-1)$.
Let $Q_{\rm tail}=(I-\Pi)Q(I-\Pi)$. Spectral invariance gives
\[
 v=G^{1/2}Q_{\rm tail}(I-\Pi)G^{1/2}\one.
\]
The bounds $\norm{Q_{\rm tail}}_{\op}\le C/\sqrt{np}$,
$\norm{(I-\Pi)G^{1/2}\one}_2\le C\sqrt n$ and
$\sum_iG_i\le Cn^2p$ imply $\sum_i|v_i|\le Cn$.
Also
\[
 \sum_i|\widehat A_{ii}|\le\sum_iG_i\Pi_{ii}
 \le2\tau K+\sum_i\frac{(G_i-\tau)^2}{G_i}\le Cn.
\]
Consequently $\sum_i|(R_1)_{ii}|\le Cn$; hollowing has total
absolute cost at most $2\sum_i|(R_1)_{ii}|$.
Since $|\widehat A_{ij}|\le a_ia_j$, \cref{restore:fit-wrong-mass} bounds
its total absolute mass by $Cn^2p$. The rank-one degree correction
has absolute mass $\sum_i|v_i|$, proving
\eqref{restore:fit-global-absolute}.

Let $E$ be the incorrectly assigned set. Each nonzero row of
$\widehat Z-Z$ has squared norm two, so the block-space residual and
\cref{restore:fit-wrong-mass} give
\[
 \norm{(I-\Pi)G^{1/2}\widehat Z}_F^2
 \le 2\norm{(I-\Pi)G^{1/2}Z}_F^2
       +2\norm{G^{1/2}(\widehat Z-Z)}_F^2
 \le C(n+np|E|).
\]
The exact identity
\[
 \widehat Z^\top(\widehat A-A)\widehat Z
 =-\{(I-\Pi)G^{1/2}\widehat Z\}^\top Q_{\rm tail}
       \{(I-\Pi)G^{1/2}\widehat Z\}
\]
therefore has norm at most $C(n+np|E|)/\sqrt{np}$.
Moving the incorrectly assigned vertices changes any raw block total
by at most $C\sum_E D_i\le C(n+np|E|)$.
Together with \eqref{restore:fit-block-event}, this bounds the error
for $\widehat A$ by $C(n+np|E|)$.

The degree correction changes a block total by
$[\widehat Z^\top v]_k n_\ell(\widehat z)/n$.
Using spectral invariance again gives
\[
 \norm{\widehat Z^\top v}_2
 \le \frac{C\sqrt{n(n+np|E|)}}{\sqrt{np}}
 \le C(n+np|E|)
\]
for all sufficiently large $np$. The global $O(n)$ hollowing cost
has already been proved. Finally, $|E|=n\Mis(\widehat z,z)$ gives
\eqref{restore:fit-total-accuracy}.
\end{proof}

For a tested vertex $i$, the local regularity condition is
\begin{equation}\label{restore:fit-local-row}
 D_i\le L_Dnp,\qquad
 a_i=\sqrt{G_i}\norm{U_{i\cdot}}_2\le L_U\sqrt p,
\end{equation}
where $L_D,L_U$ are fixed. It is imposed only at the tested target.

\begin{lemma}[Local absolute mass]
\label{restore:fit-local-mass}
On \eqref{restore:fit-local-row},
\begin{equation}\label{restore:fit-local-mass-bounds}
 \sum_j|\widehat A^{\mathrm{c}}_{ij}|\le Cnp,\qquad
 \sum_{j\in E}|\widehat A^{\mathrm{c}}_{ij}|\le C(1+p|E|)
       \quad\text{for every }E\subseteq[n].
\end{equation}
For the incorrectly assigned set the second bound is
$C\{1+np\Mis(\widehat z,z)\}$.
\end{lemma}

\begin{proof}
The inequality $|\widehat A_{ij}|\le a_ia_j$, the local bound on $a_i$ and
\cref{restore:fit-wrong-mass} prove both claims for $\widehat A$.
They also give $|v_i|\le D_i+\sum_j|\widehat A_{ij}|\le Cnp$.
The rank-one correction adds at most $Cnp$ to the full row mass and
$Cp|E|$ to its mass over $E$. Finally,
$|(R_1)_{ii}|\le a_i^2+|v_i|/n\le Cp$; hollowing changes either
bound by at most $2|(R_1)_{ii}|$. This proves the assertion.
\end{proof}

\begin{proposition}[Parameter and score errors in terms of misclassification]
\label{restore:fit-ordinary}
\label{restore:fit-score}
There is a fixed $c>0$ such that, uniformly over current labels with
$\Mis(\widehat z,z)\le c$, the symmetric Bernoulli block fit is unique
and interior for all sufficiently large $n$, and
\begin{equation}\label{restore:fit-ordinary-parameter}
 \max_{k,\ell}|\widehat P_{k\ell}-P_{k\ell}|
 \le C\left\{\frac1n+p\Mis(\widehat z,z)\right\}.
\end{equation}
At every target satisfying \eqref{restore:fit-local-row}, the fitted
log-likelihood contrast obeys
\begin{equation}\label{restore:fit-ordinary-score}
 \begin{aligned}
 &\sum_{j\ne i}\left\{
 \widehat A^{\mathrm{c}}_{ij}\log\frac{\widehat P_{k\widehat z_j}}
                              {\widehat P_{\ell\widehat z_j}}
 +(1-\widehat A^{\mathrm{c}}_{ij})\log\frac{1-\widehat P_{k\widehat z_j}}
                                   {1-\widehat P_{\ell\widehat z_j}}
 \right\}\\
 &\qquad=[\widehat A^{\mathrm{c}} w_{k\ell}]_i
       +\sum_{j\ne i}\log\frac{1-P_{kz_j}}{1-P_{\ell z_j}}
       +O\{1+np\Mis(\widehat z,z)\}.
 \end{aligned}
\end{equation}
Here $w_{k\ell}$ is the true log-odds vector in
\eqref{eq:sbm-weights}. The constants are uniform over the current
labels. In particular, a misclassification rate of order $1/(np)$
gives parameter error $O(1/n)$ and a bounded score remainder.
\end{proposition}

\begin{proof}
For one unordered fitted block pair, write $E_{k\ell}$ for its ordered
reconstructed total and $N_{k\ell}$ for the number of ordered
 off-diagonal pairs. There are two directions when $k\ne\ell$ and
one ordered within-group total when $k=\ell$. Consequently
\cref{restore:fit-block-totals} and the exact pair counts give
\[
 \begin{aligned}
 E_{k\ell}
 &=(2-\ind{k=\ell})n_kn_\ell P_{k\ell}
        +O\bigl(n\{1+np\Mis(\widehat z,z)\}\bigr),\\
 N_{k\ell}
 &=(2-\ind{k=\ell})n_k(\widehat z)n_\ell(\widehat z)+O(n).
 \end{aligned}
\]
For the common label matching,
$|n_k(\widehat z)-n_k|\le n\Mis(\widehat z,z)$.
Choose the fixed $c$ small enough that every fitted group has size
at least $n/(2\beta K)$. Thus $N_{k\ell}\asymp n^2$, and
\[
 \left|\frac{E_{k\ell}}{N_{k\ell}}-P_{k\ell}\right|
 \le C\left\{\frac1n+p\Mis(\widehat z,z)\right\}.
\]
After decreasing $c$ if necessary, this error is a sufficiently small
fixed multiple of $p$ for all large $np$. The ratio is then positive,
comparable to $p$, and strictly inside the algorithm's clipping
interval. This uses the stipulated sufficiently small $c_0$ and
$p\to0$, a consequence of $np^2\le C_1$.
The criterion $E_{k\ell}\log q+(N_{k\ell}-E_{k\ell})\log(1-q)$
has the unique interior maximizer $q=E_{k\ell}/N_{k\ell}$,
proving \eqref{restore:fit-ordinary-parameter}.

On the positive comparable-probability class, the resulting log-odds
coefficient error is at most
$C\{(np)^{-1}+\Mis(\widehat z,z)\}$.
Multiplication by the absolute row mass $Cnp$ costs at most
$C\{1+np\Mis(\widehat z,z)\}$.
Changing the incorrectly assigned neighbor labels alters bounded
coefficients on $E$, at cost at most $C(1+p|E|)$ by
\cref{restore:fit-local-mass}.
For the nonedge intercept, bounded derivatives of $\log(1-q)$ give
parameter cost at most $C\{1+np\Mis(\widehat z,z)\}$;
wrong neighbor labels add at most $Cp|E|$.
As $|E|=n\Mis(\widehat z,z)$, these bounds prove
\eqref{restore:fit-ordinary-score}. They hold simultaneously over
all stated input labels; no independence of the fitted quantities
from the tested row is used.
\end{proof}

\begin{lemma}[Uniform misclassification control for finite hard refits]
\label{restore:fit-invariant-basin}
\label{restore:fit-invariant}
There are fixed constants $c>0$ and $C<\infty$ such that, on the
aggregate event and for all sufficiently large $n$, every input
satisfying $\Mis(\widehat z,z)\le c$ yields, after one block fit and
Bernoulli assignment,
\[
 \Mis\!\left(\mathcal L(\widehat A^{\mathrm{c}},\widehat P,\widehat z),z\right)
       \le\frac{C}{np}.
\]
Starting from the prescribed initializer, all iterates satisfy a
single bound $C/(np)$ under the same label matching.
For this initialized sequence, at every target satisfying
\eqref{restore:fit-local-row}, an error at any updated iterate implies
\begin{equation}\label{restore:fit-canonical-error}
 [\widehat A^{\mathrm{c}} w_{z_i\ell}]_i+
 \sum_{j\ne i}\log\frac{1-P_{z_i z_j}}{1-P_{\ell z_j}}\le C
 \quad\text{for some }\ell\ne z_i.
\end{equation}
The threshold is the same for all such iterates and for any finite,
possibly data-dependent stopping rule after at least one update.
\end{lemma}

\begin{proof}
Choose large fixed $L_D,R$, with $L_D$ larger than twice the uniform
upper bounds on $\sum_jP_{z_i z_j}/(np)$ and $\tau/(np)$.
The degree energy bounds $|\{i:D_i>L_Dnp\}|\le C/p$.
Also $\norm{V_{i\cdot}}_2\le C/\sqrt n$, so
\eqref{restore:fit-basis} gives
$|\{i:\sqrt n\norm{U_{i\cdot}}_2>R\}|\le C/p$.
Every other target satisfies \eqref{restore:fit-local-row} for a
fixed $L_U$, because $G_i\le Cnp$ there. These counting exceptions
are independent of the input labels.

Best rank-$K$ approximation gives
$\norm{Q_K-G^{-1/2}ZPZ^\top G^{-1/2}}_{\op}\le C/\sqrt{np}$,
whence
\[
 \norm{G^{-1/2}(\widehat A-ZPZ^\top)Z}_F^2
 \le \frac{C}{np}\norm{G^{1/2}Z}_F^2\le Cn.
\]
The earlier expression for $v$ implies
$\norm{G^{-1/2}v}_2^2\le C/p$.
Since $\norm{\one^\top Z/n}_2\le1$, the degree correction
preserves the bound:
\begin{equation}\label{restore:fit-basin-action}
 \norm{G^{-1/2}(R_1-ZPZ^\top)Z}_F^2\le Cn.
\end{equation}
Hollowing changes the action of a bounded block-constant vector
at row $i$ by at most $C|(R_1)_{ii}|$.
Because $\sum_i|(R_1)_{ii}|\le Cn$, it changes that action by more
than a prescribed small multiple of $np$ on at most $C/p$ rows.
This auxiliary set is used only for counting misclassifications;
it is not charged as a probability exception in the sharp risk proof.

For $k=z_i$, the true population contrast is
\[
 [ZPZ^\top w_{k\ell}]_i+
       \sum_{j\ne i}\log\frac{1-P_{kz_j}}{1-P_{\ell z_j}}.
\]
It is a sum of incident Bernoulli Kullback--Leibler divergences plus
an $O(p)$ completed-diagonal term. Row separation and community-size
bounds give a lower bound $c_{\rm gap}np$.
At a regular target, \cref{restore:fit-ordinary} bounds the fitted
remainder by $C\{1+np\Mis(\widehat z,z)\}$, with $C$ independent
of the input labels. First choose the fixed $c$ in the hypothesis
sufficiently small compared with this population margin, then take
$np$ sufficiently large. The remainder is at most $c_{\rm gap}np/4$.
Outside the auxiliary hollowing set, a new error therefore requires
\[
 \norm{[(R_1-ZPZ^\top)Z]_{i\cdot}}_2\ge c'np
\]
for a fixed $c'>0$. Summing the squared requirements divided by
$G_i\le Cnp$ and using \eqref{restore:fit-basin-action} bounds the
number of such errors by $C/p$. Including the counting exceptions
and dividing by $n$ proves the asserted output rate. Its constant
does not depend on the input misclassification rate within the
specified small interval.

The initializer has rate at most $C/(np)$ by
\cref{restore:fit-initializer}. Choose a single fixed upper constant
for this initializer bound and the one-update output bound, then
choose $n$ large enough that its ratio to $np$ is below the fixed
$c$ just chosen. Induction now controls every iterate without
introducing another radius parameter. The estimates use the same
true community numbering at every step, so the initial common
matching is retained.

Along this sequence $np\Mis(\widehat z,z)$ is uniformly bounded.
Equation~\eqref{restore:fit-ordinary-score} thus has a single bounded
remainder for all input iterates. A wrong assignment implies
\eqref{restore:fit-canonical-error}. The same event covers finite
adaptive stopping, with no probability union over iterations.
\end{proof}
